% !Mode:: "TeX:UTF-8"

\chapter{视觉文本描述基础(测试)}
视觉文本描述是跨机器视觉和自然语言处理领域的一个挑战性课题。在描述视觉内容的过程中，通过机器视觉中的特征提取模块，将输入图像或视频转化为计算机所能理解的形式，再通过自然语言处理将视觉内容转化为自然语言的形式。在转化过程中，需要确保描述内容与视觉内容的一致性(即语义对齐)。本章介绍了视觉文本描述中常用的基本概念、理论，主要包括视觉特征提取、自然语言生成、常见的语义对齐方法、损失函数、通用标准测试集以及评价指标。从而为后续章节的算法研究奠定理论基础。

\section{视觉特征提取(测试)}
视觉特征提取是视觉文本描述中至关重要的一步。心理学的研究表明，人们可以利用整幅图像内容理解视觉内容，也可以利用图像中物品的种类和关系对视觉内容进行理解。全局特征包含完整的图像信息，但也会含有大量无关信息，比如背景信息、非重要物品信息，同时，全局特征更容易受到光照、遮挡、噪声和多视角等的影响，导致泛化能力较弱。与之相反，局部的物品特征减少了无关信息的干扰，但往往缺少物品间的空间关系信息。场景图特征在局部特征的基础上添加了物品间的关系信息。时序特征则反应了视频中的时序信息。下面将分别介绍全局特征、局部特征、场景图特征以及时序特征。

\subsection{图像全局特征(测试)}
早期的图像文本描述常采用手工特征，例如：尺度不变特征变换(Scale-invariant feature transform, SIFT)[121]、方向梯度直方图(Histogram of Oriented Gradient, HOG)[122]、颜色直方图特征、Gabor特征和空间包络特征(Gist)[122]等。

由于手工特征准确率较低且缺乏扩展能力，当前的图像文本描述方法常通过基于深度学习的方法提取全局特征。基于深度学习的特征提取模型，是将原始图像的像素值输入模型中进行学习，生成隐含图像语义的特征。该方法最大的优点在于具有鲁棒性和普适性，只需要输入图像和对应标签，模型会自动进行学习，中间无需人为干扰，提取到的特征具有更强大的表达能力，是手工特征所无法达到的。

\subsection{图像局部特征(测试)}
早期的图像文本描述常采用手工特征，例如：尺度不变特征变换(Scale-invariant feature transform, SIFT)[121]、方向梯度直方图(Histogram of Oriented Gradient, HOG)[122]、颜色直方图特征、Gabor特征和空间包络特征(Gist)[122]等。

由于手工特征准确率较低且缺乏扩展能力，当前的图像文本描述方法常通过基于深度学习的方法提取全局特征。基于深度学习的特征提取模型，是将原始图像的像素值输入模型中进行学习，生成隐含图像语义的特征。该方法最大的优点在于具有鲁棒性和普适性，只需要输入图像和对应标签，模型会自动进行学习，中间无需人为干扰，提取到的特征具有更强大的表达能力，是手工特征所无法达到的。

\section{度量学习相关理论(测试)}

\section{数据集及评价指标(测试)}

\section{本章小结(测试)}